% !TeX spellcheck = en_GB
% Sample Project Report – Federated Learning with FMI Weather Data

\documentclass[9pt]{article}
\usepackage{spconf,amsmath,bm,graphicx}
\usepackage[a4paper, margin=1in]{geometry}
\usepackage[dvipsnames]{xcolor}
\usepackage{hyperref, cleveref, tcolorbox}
\usepackage{algorithm, algorithmic}
\usepackage{amsfonts,amssymb,amsbsy}
\usepackage{subcaption, adjustbox}

\DeclareMathOperator*{\argmax}{arg\,max}
\DeclareMathOperator*{\argmin}{arg\,min}

\title{Personalized Federated Learning for Temperature Prediction Using FMI Weather Stations}
\name{Fistname Lastname}
\address{Aalto University}

\begin{document}
	\maketitle
	
	\begin{abstract}
		We study a federated learning (FL) application for short-term temperature prediction
		using weather observations from the Finnish Meteorological Institute (FMI).
		Each weather station is modeled as a node in an FL network and learns a local regression
		model from its own measurements.
		To exploit spatial correlations between nearby stations while preserving data locality,
		we employ a generalized total variation minimization (GTVMin) approach that promotes
		similarity between models of neighboring nodes.
		We implement a message-passing algorithm to solve the resulting optimization problem
		and evaluate its performance on real FMI data.
		Our results show that the proposed FL method improves prediction accuracy compared
		to purely local training, while allowing for location-specific personalization.
	\end{abstract}
	
	\textbf{Keywords:} federated learning, weather prediction, FMI data, personalized models, GTVMin
	
	\section{Introduction}
	\label{sec:intro}
	
	Accurate short-term weather prediction is essential for applications such as agriculture,
	energy management, and transportation.
	In practice, weather data are collected by geographically distributed stations that operate
	independently and may not be able or willing to share raw measurements.
	This setting naturally motivates the use of federated learning (FL), where models are trained
	collaboratively without centralizing data.
	
	In this project, we consider temperature prediction using data from the Finnish Meteorological
	Institute (FMI).
	Each weather station collects local observations such as temperature, humidity, and wind speed.
	While nearby stations are expected to exhibit similar weather patterns, local effects
	(e.g., coastal proximity or elevation) lead to station-specific behavior.
	This makes personalized FL methods particularly suitable for the task.
	
	Previous work on FL has focused on global models trained via parameter averaging.
	More recent approaches model FL systems as networks and incorporate graph-based regularization
	to balance collaboration and personalization.
	In this report, we adopt the GTVMin framework to couple local models through an empirical graph
	constructed from station locations.
	
	The remainder of this report is structured as follows.
	Section~\ref{sec:pf} formulates the FL problem.
	Section~\ref{sec:methods} describes the proposed method and algorithm.
	Section~\ref{sec:experiments} presents numerical results.
	Section~\ref{sec:conclusion} concludes the report.
	
	\section{Problem Formulation}
	\label{sec:pf}
	
	We model the application as an FL network with $N$ nodes.
	
	{\bf Nodes.}
	Each node represents one FMI weather station. 
	The node holds historical observations collected at that station only.
	
	{\bf Local models.}
	At node $i$, we train a linear regression model
	\[
	\hat{y}^{(i)} = \mathbf{w}^{(i)\top} \mathbf{x},
	\]
	where $\mathbf{x}$ contains features such as past temperature, humidity, and wind speed,
	and $\hat{y}^{(i)}$ is the predicted temperature at the next time step.
	
   {\bf Loss functions.}
	Each node uses a local mean squared error (MSE) loss
	\[
	f_i(\mathbf{w}^{(i)}) =
	\frac{1}{|\mathcal{D}_i|}
	\sum_{(\mathbf{x},y)\in \mathcal{D}_i}
	\bigl(y - \mathbf{w}^{(i)\top}\mathbf{x}\bigr)^2,
	\]
	where $\mathcal{D}_i$ denotes the local dataset at node $i$.
	
	{\bf Edges.}
	Edges between nodes are defined based on geographical proximity.
	We connect each station to its $k$ nearest neighboring stations using Euclidean distance
	between coordinates.
	Edge weights are chosen as a decreasing function of distance, assigning larger weights
	to closer stations.
	
	\section{Methods}
	\label{sec:methods}
	
	We apply a GTVMin-based formulation to jointly learn personalized local models.
	The optimization problem is given by
	\[
	\min_{\{\mathbf{w}^{(i)}\}}
	\sum_{i=1}^N f_i(\mathbf{w}^{(i)})
	+ \lambda \sum_{(i,i')\in \mathcal{E}}
	\phi\!\left(\mathbf{w}^{(i)} - \mathbf{w}^{(i')}\right),
	\]
	where $\mathcal{E}$ denotes the set of edges in the FL network.
	
	{\bf Variation measure.}
	We choose the squared $\ell_2$ norm
	\[
	\phi(\mathbf{w}^{(i)} - \mathbf{w}^{(i')})
	= \|\mathbf{w}^{(i)} - \mathbf{w}^{(i')}\|_2^2,
	\]
	which encourages smooth variation of model parameters across neighboring stations.
	
	{\bf FL algorithm.}
	The problem is solved using a message-passing algorithm derived from gradient-based
	optimization.
	Each node alternates between a local gradient step on its loss function and an
	aggregation step where it exchanges parameter information with neighboring nodes.
	This implementation avoids sharing raw weather data and only communicates model updates.
	
	\section{Numerical Experiments}
	\label{sec:experiments}
	
	{\bf Data.}
	We use publicly available FMI open data, including hourly measurements of temperature,
	humidity, and wind speed from multiple weather stations across Finland.
	The dataset is split into training, validation, and test sets in chronological order.
	
	{\bf Model selection and validation.}
	The regularization parameter $\lambda$ and the number of neighbors $k$ are selected using
	validation data.
	We monitor training and validation losses to detect overfitting.
	
	{\bf Results.}
	Compared to purely local training, the proposed FL method achieves lower test MSE at most
	stations.
	Stations with limited local data benefit most from collaboration, while the method still
	preserves station-specific differences in coastal and inland regions.
	
	\section{Conclusion}
	\label{sec:conclusion}
	
	We investigated a personalized FL approach for temperature prediction using FMI weather data.
	By modeling weather stations as nodes in an FL network and applying GTVMin-based regularization,
	we successfully leveraged spatial correlations while maintaining local autonomy.
	
	The results demonstrate improved predictive performance compared to local-only models.
	Limitations of this study include the use of simple linear models and static graph
	construction.
	Future work could explore nonlinear models, adaptive graphs, and additional meteorological
	variables.
	
	\begin{thebibliography}{9}
		\bibitem{Jung2025}
		A. Jung,
		\textit{Federated Learning: From Theory to Practice},
		Aalto, 2025.
		Available:
		\url{https://github.com/alexjungaalto/FederatedLearning/blob/main/material/FLBook.pdf}.
		
		\bibitem{Jung2022}
		A. Jung,
		\textit{Machine Learning: The Basics},
		Springer, 2022.
	\end{thebibliography}
	
\end{document}
